%% INTRO — Draft B: Geometry-Aware Action Tokenization framing
\section{Introduction}
\label{sec:intro}

Discrete action tokenization has become a cornerstone of modern VLA models.
\fast{}~\cite{fast2025} achieves $27.8\times$ compression via DCT+BPE,
enabling dexterous tasks where per-dimension binning fails.
\oat{}~\cite{oat2026} introduces ordered token spaces via register encoders and
nested dropout, enabling prefix-decodable autoregressive policies.
These advances have driven VLA accuracy to $94$--$98\%$ on LIBERO~\cite{libero2023}.

\textbf{Yet all current action tokenizers share a critical blind spot: they are
geometry-unaware.}
\fast{} tokenizes raw world-frame trajectories; \oat{} tokenizes raw action space.
A trajectory translated $10\,\text{cm}$ produces an entirely different token stream.
Raw action tokens change $99.1\%$ of their identity under random rigid-body transforms ---
spatial robustness cannot emerge from such representations by construction,
which explains the 14--18\% performance drops observed across all VLA models
under spatial distribution shift.

\textbf{We show that the fix is to tokenize in $\se{3}$-invariant space.}
\dhb{} encodes trajectories via the Double-Reflection invariant frame,
separating motion shape from global pose.
Magnitude tokens change $0.0\%$ under any rigid-body transform.
The invariant token stream enables geometric adaptation at $\leq\!7\,\text{ms}$:
re-ground the encoded shape to any new target pose without retraining or re-encoding.
\fast{} can compress actions; \dhb{} can generalize across poses.

\textit{FAST compresses in raw space and cannot generalize spatially.
DHB-XR tokenizes in invariant space and adapts geometrically at 7\,ms.}

\paragraph{Contributions.}
\begin{enumerate}
\item \textbf{Geometry-aware invariant tokenization.}
First action tokenizer with provable SE(3)-invariance: $0.0\%$ magnitude token change
vs $99.1\%$ raw, with 8-token OAT-style compression and prefix decodability.

\item \textbf{Spatial generalization.}
$0.0\,\text{mm}$ goal error at $10\,\text{cm}$ perturbation vs $100.0\,\text{mm}$ naive,
validated end-to-end in simulation.

\item \textbf{StarVLA ablation head.}
\texttt{DHBActionHeader} as a drop-in StarVLA module: clean ablation of FAST head
vs MLP head vs \dhb{} head on the same Qwen2.5-VL backbone.

\item \textbf{Sample efficiency.}
$N{=}1$ invariant demo covers any pose neighborhood;
same generalization as $N{=}20$ raw demos.
\end{enumerate}
